\chapter{Introduction}
Sign languages are the primary languages of Deaf communities, and they
are full natural languages rather than manual renderings of the
surrounding spoken language. The World Federation of the Deaf estimates
that more than 70 million deaf people belong to signing communities,
using over 200 national sign languages~\cite{wfd2026faq}.

Automatic processing of these languages has been studied extensively, and
recent surveys document an active literature spanning recognition,
translation and
generation~\cite{bragg2019sign, koller2020quantitative,
nunezmarcos2023survey, decoster2024machine}. Widely usable tools,
however, do not yet exist, and reported progress rests substantially on
having seen the signer before: recent translation systems lose most of
their accuracy when re-evaluated on signers unseen in
training~\cite{artiagaRethinkingSignLanguage2025}. Closing that gap needs
representations that generalise beyond the signers and recording
conditions any one corpus contains, which is the prerequisite this thesis
works on.

\emph{Sign language recognition} (SLR) maps video of signing to
linguistic units. It is conventionally split into \emph{isolated} SLR,
where a clip contains a single sign and the model predicts a \emph{gloss}
(a written label for a sign, by convention in capitals) from a fixed
vocabulary, and \emph{continuous} SLR, where a sentence-length recording
is transcribed as a gloss sequence; both are distinct from translation,
which targets spoken-language text. This thesis works at the isolated
level on one corpus and the sentence level on the other, evaluating both
by retrieval (\S\ref{sec:eval-protocols}).

SLR faces two distinct problems that this
thesis addresses together.
The first is \emph{data scarcity}, whose scale is easy to understate. The
largest sign language video corpus holds roughly $3{,}200$ hours, and
covers only about twenty-five of the more than three hundred sign
languages in use worldwide~\cite{tanzerYouTubeSL25}; a single speech
recognition model of the same period was trained on $680{,}000$ hours of
audio~\cite{radfordRobustSpeechRecognition2022}. Annotated data is
scarcer again: the largest corpora are weakly supervised harvests whose
text is scraped or machine-transcribed rather than
glossed~\cite{tanzerYouTubeSL25, liUniSignUnifiedSign2025}, and glossing
proper is slow expert work, about an hour per ninety seconds of video for
How2Sign~\cite{duarteHow2Sign2021}, with formats that are not
standardised across
corpora~\cite{desistoChallengesSignLanguage2022, DataProblem}. 

The second is \emph{appearance bias}: a video model can score well on a benchmark by
exploiting the static, incidental visual properties of whoever happened to
be recorded (clothing, skin tone, body shape, background) rather than focusing exclusively on
the motion that carries the linguistic content. In the action-recognition
literature this second phenomenon is well
documented~\cite{leiRevealingSingleFrame2023, ilicAppearanceFreeAction2022,
liMitigatingEvaluatingStatic2023a, zhangDontJudgeLook2024}; action
recognition is also the domain closest to ours (SLR inherits its video
backbones and training recipes from it), so we treat its bias findings
and interventions as directly relevant, and review them in
\S\ref{sec:related-appearance-bias}. The two
problems interact (fewer recorded signers make appearance shortcuts
easier to exploit), but they are separate, and this thesis uses one
mechanism against both: synthetically generated signing adds data where
data is scarce, and its controlled appearance variation attacks the bias
directly.

For sign language this is not an abstract concern. Signing is performed by
people, and the appearance dimensions a model may latch onto (skin tone,
gender, clothing, recording environment) correlate directly
with demographic groups~\cite{heldrethWhichSkinTone2024,
schumannConsensusSubjectivitySkin2023}. A sign language recognition system that implicitly
keys on appearance will degrade unpredictably for signers who look unlike its
training population. Signer variation itself is a long-standing focus of
SLR research, and signer-independent evaluation is an established
protocol (\S\ref{sec:related-signer-independent} reviews the remedies
proposed). Yet while appearance bias is actively studied in general
action recognition, it has received little direct attention in the sign
language domain.

The field's dominant response to signer variation has been indirect: move
from RGB input to pose-based input. Pose representations discard pixels
(and with them, most appearance information), and a visible trend in recent
benchmarks and challenges pushes towards signer-independent evaluation and
pose-based methods~\cite{jangSigningOutsideStudio2022,
luqmanSignEval2025Challenge2025a, jiangSignCLIPConnectingText2024}. The
SignEval 2025 challenge makes the trend explicit: in its signer-independent
recognition task, every submitted system was pose-based, by design,
since the organisers distributed only MediaPipe-derived pose features
extracted from the underlying (RGB-recorded) Isharah videos, citing signer
overfitting~\cite{luqmanSignEval2025Challenge2025a,
alyamiIsharahLargeScaleMultiScene2026b}: for that task, RGB-based
approaches were excluded from the outset rather than tried and outperformed.

But pose is not a free lunch. Pose extractors are themselves learned
RGB models, with demographic and mechanical failure modes of their own
(\S\ref{sec:related-work-modality}). Meanwhile recent self-supervised
RGB representations report stronger sign representations than pose-based
alternatives~\cite{wongSignRepEnhancingSelfSupervised2025}. If RGB has the higher ceiling, then abandoning it to escape
appearance bias trades accuracy for robustness. The alternative this thesis
explores is to keep RGB and attack the bias directly with \emph{synthetic
appearance counterfactuals}: renderings or edits of the same signing motion
under different appearances, so that appearance becomes a factor the model
can be explicitly trained, or shown, to ignore.

The approach we set out to build is the one that guarantees the
counterfactual property by construction: record signing with motion
capture, retarget the captured motion to synthetic avatars, and render
them. Every render then shows exactly the same signing motion under a
freely chosen appearance, so motion preservation is a property of the
pipeline rather than an assumption; we call these \emph{motion-locked
re-renders}. But this route cannot be
applied to an existing benchmark where videos are plain RGB with no
motion capture attached, so the only way to use it at benchmark scale would
be to record a new mocap corpus of comparable size.

We therefore approximate it. Modern generative image editors can rewrite
the appearance of an existing video frame by frame: we edit
real benchmark videos (changing clothing, adding glasses, altering skin
tone, swapping in a different signer) while leaving the rest of each
frame alone. This recovers the appearance-variation axis of the
re-renders at a cost that scales to existing benchmarks. It is an
approximation, not a substitute: the edits are per-frame appearance
rewrites, so the signing motion is preserved only insofar as the editor
leaves it untouched, not by construction.

\begin{figure*}[t]
\centering
\includegraphics[width=\textwidth]{method_overview.pdf}
\caption[Method overview: appearance counterfactuals in contrastive
training]{\textbf{Method overview.} Setting A is the high-diversity
regime, Setting B the low-diversity one. Objectives and loss weights are defined in
\S\ref{sec:method-objectives}.}
\label{fig:method_overview}
\end{figure*}

In this thesis we therefore study two dataset regimes that differ in
appearance diversity, both with signer-disjoint evaluation. The first is
a low-diversity setting in which motion capture is available: NGT744, a
small, self-recorded Dutch Sign Language (NGT) corpus recorded from four
signers in a single studio and evaluated in another, so that both signer
identity and capture conditions are near-constant during training. There
the motion-locked pipeline runs in full, and we compare the two
counterfactual routes head to head, avatar re-renders against generative
edits of the same videos, under matched conditions. The second is a
high-diversity setting in which only RGB video is available: ASL Citizen,
whose 35 training signers record themselves on their own webcams, so that
identity, lighting, framing and background all vary. There only the
generative route applies, and we ask whether synthetic appearance data
still helps. NGT744 is additionally low-resource in absolute size; ASL
Citizen is not large by general-vision standards, but it is diverse. The
two settings turn out to give different answers, and that contrast, where
synthetic appearance variation earns its cost and where it does not, is
the core result of the thesis.

This thesis asks whether appearance bias is measurable in modern sign
language representations, and whether synthetic appearance
counterfactuals can reduce it; Figure~\ref{fig:method_overview}
summarises how they enter training. Concretely:

\begin{description}
\item[RQ1 (Diagnosis).] Do sign language video representations, namely pose-based
  (MediaPipe), RGB transformer
  (Logos/MViTv2 \cite{ovodovLogosWellTemperedPretrain2025}), RGB CNN (I3D),
  and contrastively text-aligned (SignCLIP \cite{jiangSignCLIPConnectingText2024}),
  encode signer identity beyond sign content, and to what degree?

\item[RQ2 (Low-diversity regime).] Do synthetic appearance
  counterfactuals improve retrieval where training diversity is minimal
  (the NGT744 corpus, four signers in a single studio)? We ask this of
  each synthesis route on its own before comparing them:
  \begin{description}
  \item[RQ2a.] Do motion-locked re-renders (avatar renders driven by real
    motion capture) help when added as extra views under a supervised
    contrastive objective~\cite{khosla2020supervised}?
  \item[RQ2b.] Do generative appearance edits help in the same role?
  \item[RQ2c.] At matched strength, content and viewpoint, how do the two
    routes compare?
  \end{description}

\item[RQ3 (High-diversity regime).] In a high-diversity setting (ASL
  Citizen), do generative appearance edits
  (clothing, glasses, skin tone, signer swap), our
  scalable approximation of motion-locked synthesis, improve
  retrieval, either as counterfactual supervision when fine-tuning the backbone itself or as training data for a contrastive embedding head?

\item[RQ4 (Interpretation).] What do the RQ3 interventions do
  to the embedding space itself: where do the synthetic counterfactuals
  land relative to their source videos and to a genuine signer change
  (do they stay local, and do they occupy a separable ``synthetic'' region),
  and does any training objective actually reduce the signer identity
  encoded in the features, measured both by distance geometry and by
  linear decodability?
% Answerable by: \S\ref{sec:results-l2-raw} + \S\ref{sec:results-ft-probes}
% (distance ratio + linear decoding, incl. the dissociation, lines
% 1309-1380 and 2002-2108) and \S\ref{sec:results-aug-geometry}
% (displacement locality, sim-to-real separability, de-shift collapse,
% lines 2110-2323). "genuine signer change" echoes the n.disp calibration
% (lines 2141-2151); "separable synthetic region" echoes measurement 2
% (lines 2158-2174); "distance geometry ... linear decodability" echoes the
% two probes (lines 2008-2021).
\end{description}



